% === quiz_solutions.tex ===
\section*{Answer Key: Vector Calculus and Integral Theorems Activity}

\begin{enumerate}
    \item Let
          \[
              f(x,y)=x^2+xy+y^2.
          \]
          First compute the gradient:
          \[
              \nabla f=\left\langle 2x+y,\; x+2y\right\rangle
          \]
          At $(1,2)$:
          \[
              \nabla f(1,2)=\langle 4,\; 5\rangle
          \]
          The magnitude of $\langle 3,4\rangle$ is
          \[
              \sqrt{3^2+4^2}=5
          \]
          so the unit vector is
          \[
              \mathbf{u}=\left\langle \frac35,\; \frac45\right\rangle
          \]
          Hence the directional derivative is
          \[
              D_{\mathbf{u}}f=\nabla f\cdot \mathbf{u}
              =\langle 4,5\rangle\cdot \left\langle \frac35,\frac45\right\rangle
              =\frac{12}{5}+\frac{20}{5}
              =\frac{32}{5}
          \]
          \[
              \boxed{\textcolor{red}{\frac{32}{5}}}
          \]

    \item We are given
          \[
              \mathbf{F}(x,y)=\langle y,\; x\rangle,\qquad
              \mathbf{r}(t)=\langle t,\; t^2\rangle,\quad 0\le t\le 1
          \]
          Differentiate the curve:
          \[
              \mathbf{r}'(t)=\langle 1,\; 2t\rangle
          \]
          Evaluate the vector field along the curve:
          \[
              \mathbf{F}(\mathbf{r}(t))=\langle t^2,\; t\rangle
          \]
          Compute the dot product:
          \[
              \mathbf{F}(\mathbf{r}(t))\cdot \mathbf{r}'(t)=t^2(1)+t(2t)=3t^2
          \]
          Therefore,
          \[
              \int_C \mathbf{F}\cdot d\mathbf{r}
              =\int_0^1 3t^2\,dt
              =\left[t^3\right]_0^1
              =1
          \]
          \[
              \boxed{\textcolor{red}{1}}
          \]

    \item By Green's Theorem,
          \[
              \oint_C (P\,dx+Q\,dy)=\iint_R \left(\frac{\partial Q}{\partial x}-\frac{\partial P}{\partial y}\right)\,dA
          \]
          where
          \[
              P=x^2y,\qquad Q=xy^2
          \]
          Compute the needed partial derivatives:
          \[
              \frac{\partial Q}{\partial x}=y^2,\qquad
              \frac{\partial P}{\partial y}=x^2
          \]
          Hence,
          \[
              \oint_C (x^2y\,dx+xy^2\,dy)=\iint_R (y^2-x^2)\,dA
          \]
          The triangular region is bounded by
          \[
              0\le x\le 1,\qquad 0\le y\le 2x
          \]
          so
          \[
              \iint_R (y^2-x^2)\,dA
              =
              \int_0^1\int_0^{2x}(y^2-x^2)\,dy\,dx
          \]
          Evaluate the inner integral:
          \[
              \int_0^{2x}(y^2-x^2)\,dy
              =
              \left[\frac{y^3}{3}-x^2y\right]_0^{2x}
              =\frac{8x^3}{3}-2x^3
              =\frac{2x^3}{3}
          \]
          Then
          \[
              \int_0^1 \frac{2x^3}{3}\,dx
              =\frac23\left[\frac{x^4}{4}\right]_0^1
              =\frac{1}{6}
          \]
          \[
              \boxed{\textcolor{red}{\frac16}}
          \]

    \item By Stokes’ Theorem,
          \[
              \oint_C \mathbf{F}\cdot d\mathbf{r}
              =
              \iint_S (\nabla\times \mathbf{F})\cdot d\mathbf{S}
          \]
          where
          \[
              \mathbf{F}=\langle -y,\; x,\; z\rangle
          \]
          Compute the curl:
          \[
              \nabla\times \mathbf{F}
              =
              \begin{vmatrix}
                  \mathbf{i} & \mathbf{j} & \mathbf{k} \\
                  \partial_x & \partial_y & \partial_z \\
                  -y         & x          & z
              \end{vmatrix}
              =\langle 0,\; 0,\; 2\rangle
          \]
          The surface $S$ is the disk $x^2+y^2\le 1$ in the plane $z=1$, oriented upward, so
          \[
              \mathbf{n}=\langle 0,0,1\rangle
          \]
          and
          \[
              (\nabla\times \mathbf{F})\cdot \mathbf{n}
              =\langle 0,0,2\rangle\cdot \langle 0,0,1\rangle
              =2
          \]
          Therefore,
          \[
              \iint_S (\nabla\times \mathbf{F})\cdot d\mathbf{S}
              =\iint_S 2\,dA
              =2(\text{area of unit disk})
              =2\pi
          \]
          \[
              \boxed{\textcolor{red}{2\pi}}
          \]

    \item By the Divergence Theorem,
          \[
              \iint_S \mathbf{F}\cdot d\mathbf{S}
              =
              \iiint_V (\nabla\cdot \mathbf{F})\,dV
          \]
          where
          \[
              \mathbf{F}=\langle x^3,\; y^3,\; z^3\rangle
          \]
          Compute the divergence:
          \[
              \nabla\cdot \mathbf{F}
              =
              \frac{\partial}{\partial x}(x^3)+\frac{\partial}{\partial y}(y^3)+\frac{\partial}{\partial z}(z^3)
              =3x^2+3y^2+3z^2
          \]
          On the unit sphere, spherical coordinates are convenient:
          \[
              x^2+y^2+z^2=r^2,\qquad dV=r^2\sin\theta\,dr\,d\theta\,d\phi
          \]
          Thus,
          \[
              \iiint_V (\nabla\cdot \mathbf{F})\,dV
              =
              \int_0^{2\pi}\int_0^\pi\int_0^1 3r^2(r^2\sin\theta)\,dr\,d\theta\,d\phi
          \]
          \[
              =
              3\left(\int_0^1 r^4\,dr\right)
              \left(\int_0^\pi \sin\theta\,d\theta\right)
              \left(\int_0^{2\pi} d\phi\right)
          \]
          Evaluate each integral:
          \[
              \int_0^1 r^4\,dr=\frac15,\qquad
              \int_0^\pi \sin\theta\,d\theta=2,\qquad
              \int_0^{2\pi} d\phi=2\pi
          \]
          Therefore,
          \[
              3\cdot \frac15\cdot 2\cdot 2\pi
              =\frac{12\pi}{5}
          \]
          \[
              \boxed{\textcolor{red}{\frac{12\pi}{5}}}
          \]
\end{enumerate}